\documentclass[11pt]{article}
\usepackage[margin=1in]{geometry}
\usepackage{amsmath,amssymb,amsthm}
\newtheorem{theorem}{Theorem}
\newtheorem{lemma}{Lemma}
\newtheorem{corollary}{Corollary}
\newtheorem{definition}{Definition}
\title{V117: A W[1]-Hard Two-Chain Feedback-Linkage Barrier}
\author{NC0\_k-Avoid Laboratory}
\date{}
\begin{document}
\maketitle

\paragraph{Scope.}
This note transfers a published DAG edge-disjoint-path hardness theorem into the laboratory's gate-resource residual model. It does not prove W[1]-hardness of the full signed-MUX problem, unrestricted $NC^0_3$-Avoid, or any resolution of P versus NP.

\begin{definition}[Parallel two-chain gate linkage]
A DAG has unrestricted variable vertices and capacity-one gate resources. The requests are partitioned into two ordered chains $C_0,C_1$. Paths for all requests must be globally disjoint on gate resources; variable vertices may be reused. In the parallel promise, all requests in one chain have the same ordered terminal pair.
\end{definition}

\begin{theorem}[Two-chain gate-linkage barrier]
Parallel two-chain gate linkage on DAGs is W[1]-hard parameterized by the total number $k=|C_0|+|C_1|$ of requests.
\end{theorem}
\begin{proof}
Slivkins~\cite{Slivkins2010} proves that Edge-Disjoint Paths on DAGs is W[1]-hard in the number of requests even when the demand graph consists of two sets of parallel edges. We may assume both sets are nonempty. If one set is empty, add two fresh vertices joined by one private edge and put the forced request between them into the missing set. This preserves acyclicity and feasibility and increases the request parameter by only one.

For every source edge $e=(u,v)$ insert one fresh capacity-one gate resource $g_e$ and replace $e$ by $u\to g_e\to v$. Put the two nonempty parallel demand classes into $C_0,C_1$.

The subdivision is acyclic because contracting all $g_e$ recovers the source DAG. A source path uses $e$ exactly when its subdivision uses $g_e$. Hence source paths are pairwise edge-disjoint iff the transformed request paths are pairwise gate-resource-disjoint. The request parameter is preserved up to the optional additive-one padding, hence this is an FPT parameter-preserving reduction.
\end{proof}

\begin{definition}[Prescribed feedback closure]
For each consecutive pair of requests in a chain, add a fresh private feedback gate from the preceding request target to the next request source. A solution for each chain is required to traverse those feedback gates in the prescribed order. Let $F$ be the set of these feedback gates.
\end{definition}

\begin{theorem}[Supplied feedback-interface barrier]
If both chains are nonempty, the feedback closure has $|F|=k-2$. Deleting $F$ leaves the gate-subdivided DAG, and feasibility is equivalent to the two-chain instance. Consequently the prescribed two-route residual problem is W[1]-hard parameterized by the supplied feedback-interface size $\tau=|F|$.
\end{theorem}
\begin{proof}
A chain with $r$ requests receives $r-1$ private feedback gates. Concatenating its $r$ gate-disjoint request paths with those gates gives one gate-simple prescribed walk. Conversely, splitting such a walk at its prescribed feedback gates recovers the $r$ request paths. Summing over two nonempty chains gives $\tau=(|C_0|-1)+(|C_1|-1)=k-2$. With the optional padding above, $\tau$ differs from the original hard parameter by at most an additive constant.
\end{proof}

\begin{corollary}
Unless $\mathrm{FPT}=\mathrm{W[1]}$, the V116 variable-dimensional residual routing problem cannot be compressed to $f(\tau)N^{O(1)}$ using only the facts that deletion of a supplied $\tau$-gate interface leaves a DAG and the residual requests form two prescribed chains.
\end{corollary}

\paragraph{Non-implication.}
The reduction is at the gate-resource routing level. It does not prove that the hard instances can be realized as signed essential ternary MUX dominator stages, nor that the supplied feedback set is a minimum feedback-gate set. Those are the V118 questions.

\begin{thebibliography}{9}
\bibitem{Slivkins2010}
A.~Slivkins, \emph{Parameterized Tractability of Edge-Disjoint Paths on Directed Acyclic Graphs}, SIAM J. Discrete Math. 24(1):146--157, 2010. DOI 10.1137/070697781.
\end{thebibliography}
\end{document}
